% Part: methods
% Chapter: proofs
% Section: cant-do-it

\documentclass[../../../include/open-logic-section]{subfiles}

\begin{document}

\olfileid{mth}{prf}{cnt}

\olsection{నా వల్ల కావడం లేదు!{}}

ఏదో ఒక సమయంలో మనందరికీ వదిలేయాలనిపిస్తుంది. కానీ మీరు \emph{చేయగలరు}. మీ అధ్యాపకులు, బోధనా సహాయకులు, తోటి విద్యార్థులు మీకు తోడ్పడగలరు. సహాయం అడగండి!{} సంక్షోభం రాకుండా ఉండటానికి, వదిలేయాలనిపించినప్పుడు ఏమి చేయాలో కొన్ని సూచనలు ఇవి.

సమస్యలను విజయవంతంగా పరిష్కరించడానికి ఇలా చేయండి:
\begin{enumerate}
\item \emph{సాధ్యమైనంత ముందుగానే మొదలుపెట్టండి.} సెమిస్టర్ అంతటా పనులు ఎక్కువగా ఉంటాయి; పనిని వాయిదా వేయడం కూడా చాలామందికి కష్టమే. కాబట్టి ఇంటిపనిని ముందే మొదలుపెట్టడం ఉత్తమమైన పనులలో ఒకటి. అప్పుడు ఎక్కడైనా చిక్కుకుంటే (ఏడవడం కాకుండా) పరిష్కారం వెతకడానికి సమయం ఉంటుంది.
\item \emph{తోటి విద్యార్థులతో మాట్లాడండి.} మీరు ఒంటరివారు కాదు. తరగతిలో ఇతరులకూ కష్టాలు ఉండవచ్చు---కాని వారికి వేరే విషయాల్లో కష్టం ఉండవచ్చు. సమస్య గురించి వారితో మాట్లాడితే కొత్త దృష్టికోణం దొరికి పరిష్కార మార్గం కనిపించవచ్చు. అయితే వారి సమాధానాన్ని కేవలం కాపీ చేయవద్దు: సూచన అడగండి; లేదా ఎక్కడ ఆగిపోయారో వివరించి తరువాతి దశ గురించి అడగండి. మీకు అర్థమైన తరువాత వారికీ అలాగే తోడ్పడండి. ఇతరులకు వివరించి సహాయం చేయడం మీ అవగాహననూ పెంచుతుంది.
\item \emph{సహాయం అడగండి.} మీకు అనేక వనరులు ఉన్నాయి---అధ్యాపకులు, బోధనా సహాయకులు మీకు తోడుగా ఉంటారు; మీరు విజయం సాధించాలని \emph{కోరుకుంటారు}. సమస్యను విడమరచి చూడటానికి, ఏ దశలో కష్టపడుతున్నారో గుర్తించడానికి వారు సహాయపడగలరు.
\item \emph{విరామం తీసుకోండి.} సమస్యనే చాలా సేపు చూస్తూ ఉండటం వల్ల కూడా మీరు ఆగిపోయి \emph{ఉండవచ్చు}. కొంత విరామం తీసుకోండి; ఒక కప్పు టీ తాగండి; లేదా కొంతసేపు మరో సమస్యపై పనిచేసి, తాజా మనసుతో దీనికి తిరిగి రండి. రాత్రి విశ్రాంతి తీసుకొని మరుసటి రోజు చూడండి.
\end{enumerate}

ఈ పద్ధతులు పనిచేయాలంటే నిరూపణపై బాగా ముందుగానే పనిచేయడం మొదలుపెట్టి ఉండాలని గమనించారా? గడువు ముందురోజు తెల్లవారుజామున రెండు గంటలకు మొదలుపెడితే ఇవి అంతగా ఉపయోగపడకపోవచ్చు.

ఇది నిరాశాజనకంగా వినిపించవచ్చు; కానీ నిరూపణ సమస్యను పరిష్కరించడం చివరకు ఫలించే సవాలు. కొందరు దీన్నే వృత్తిగా చేసుకుంటారు---అంటే ఇందులో ఆస్వాదించదగినది ఏదో ఉండాలి. దాదాపు అన్నిటిలాగే సమస్యలను పరిష్కరించడానికి, నిరూపణలు చేయడానికి సాధన అవసరం. ఇది సులభంగా చేసే తోటి విద్యార్థులు కనిపించవచ్చు: వారు బహుశా ఇప్పటికే చాలా సాధన చేసి ఉంటారు. త్వరగా చేతులెత్తేయకండి.

ఒక సమస్యకు సమయం (లేదా ఓపిక) అయిపోయినా పరవాలేదు. దాని అర్థం మీకు తెలివి లేదనో, దాన్ని ఎప్పటికీ అర్థం చేసుకోలేరనో కాదు. అది ఎలా చేస్తారో అధ్యాపకుల నుంచి లేదా మరో విద్యార్థి నుంచి తెలుసుకోండి; మీరు ఎక్కడ తప్పు చేశారో, ఎక్కడ ఆగిపోయారో గుర్తించండి. అప్పుడు ఇలాంటి సమస్య మళ్లీ ఎదురైతే అదే ఇబ్బంది రాకుండా చూసుకోవచ్చు. తరువాత సమాధానం చూడకుండా మీరే చేయడానికి ప్రయత్నించండి. వచ్చేసారి ముందుగానే మొదలుపెట్టి, సహాయమూ ముందుగానే అడగండి.

\end{document}
